\documentclass[12pt]{article}
\usepackage[T1]{fontenc}
\usepackage[T1]{polski}
\usepackage[utf8]{inputenc}
\newcommand{\BibTeX}{{\sc Bib}\TeX} 
\usepackage{graphicx}
\usepackage{amsfonts}
\usepackage{float}
\usepackage{amsmath}
\usepackage{subcaption}

\setlength{\textheight}{21cm}

\title{{\bf Zadanie nr 4 - Przekształcenie Fouriera, Walsha-Hadamarda, kosinusowe i falkowe,
szybkie algorytmy.}\linebreak
Cyfrowe Przetwarzanie Sygnałów}
\author{Maciej Miazek 251589 \and Remigiusz Tomecki 251652}
\date{\today} 

\begin{document}
\clearpage\maketitle
\thispagestyle{empty}
\newpage
\setcounter{page}{1}

\section{Cel zadania}
Celem niniejszego ćwiczenia było zapoznanie się z fundamentalnymi operacjami transformacji sygnałów dyskretnych z dziedziny czasu w dziedzinę częstotliwości (i pokrewne). Wymagane było praktyczne zaimplementowanie algorytmów Dyskretnej Transformacji Fouriera (DFT) oraz jej zoptymalizowanych, szybkich wariantów z decymacją w czasie (DIT FFT) i częstotliwości (DIF FFT). Analogiczne zoptymalizowane podejścia (algorytmy szybkie w złożoności $O(N \log_2 N)$) należało wykorzystać do implementacji Transformacji Kosinusowej (DCT i FCT) oraz Walsha-Hadamarda (WHT i FWHT). Ważnym elementem ćwiczenia była również modyfikacja istniejącej aplikacji w celu obsługi sygnałów o wartościach zespolonych oraz odpowiedniej interpretacji widma amplitudowego i fazowego.

\section{Wstęp teoretyczny}
Głównym aparatem matematycznym stosowanym w cyfrowym przetwarzaniu sygnałów do analizy widmowej jest Dyskretna Transformacja Fouriera. Przekształca ona wektor $x(n)$ długości $N$ na zbiór prążków $X(m)$:
\begin{equation}
X(m) = \frac{1}{N} \sum_{n=0}^{N-1} x(n) e^{-j \frac{2\pi m}{N} n}, \quad 0 \le m \le N-1
\end{equation}
Standardowy algorytm DFT realizuje te obliczenia w złożoności czasowej $O(N^2)$. Wykorzystując symetrię i okresowość czynnika obrotu (tzw. jądra transformacji) $W_N = e^{j \frac{2\pi}{N}}$, algorytm Szybkiej Transformacji Fouriera (FFT) pozwala zredukować złożoność do $O(N \log_2 N)$ m.in. dla wektorów o długości $N=2^k$.

Transformacja Kosinusowa drugiego rodzaju (DCT-II) ma szerokie zastosowanie w metodach kompresji (np. format JPEG). Opiera się wyłącznie na liczbach rzeczywistych:
\begin{equation}
X_m = c(m) \sum_{n=0}^{N-1} x(n) \cos\left( \frac{\pi (2n+1) m}{2N} \right)
\end{equation}
Z kolei Transformacja Walsha-Hadamarda (WHT) wykorzystuje w pełni prostokątne i dyskretne funkcje bazowe złożone z wartości $1$ oraz $-1$. Macierze te definiowane są rekurencyjnie:
\begin{equation}
H_0 = [1], \quad H_m = \frac{1}{\sqrt{2}} \begin{bmatrix} H_{m-1} & H_{m-1} \\ H_{m-1} & -H_{m-1} \end{bmatrix}
\end{equation}
Szybki wariant tej transformaty (FWHT) polega wyłącznie na naprzemiennym dodawaniu i odejmowaniu elementów wektora, co powoduje brak zapotrzebowania na kosztowne układy mnożące w systemach sprzętowych.


\section{Opis implementacji}
Aplikacja została zaimplementowana w języku Python. Wykorzystane zostały biblioteki \verb|Numpy| do wykonywania szybkich operacji macierzowych oraz wektorowych, a interfejs graficzny powstał w \verb|Tkinter|.

Do głównego kodu programu wprowadzono:
\begin{itemize}
    \item \textbf{Generację złożonych sygnałów zespolonych}: Zdefiniowano operatory $S1, S2, S3$ symulujące sumę funkcji sinusoidalnych, wprowadzono również ogólną obsługę i wizualizację sygnałów jako wartości zespolonych wektora \verb|numpy.array|.
    \item \textbf{Klasę \texttt{TransformacjaSygnalu}}: W ramach tej klasy zamknięto metody realizujące odpowiednie wzory matematyczne. Bezpośrednie podejścia korzystają z \verb|numpy.dot()| w celu efektywnego mnożenia. Szybkie wersje algorytmów (DIT, DIF) zaimplementowano z wykorzystaniem struktury "motylkowej" (metody rekurencyjne bądź "in-place"), rozdzielając operacje na składowe parzyste i nieparzyste.
    \item \textbf{Moduł graficzny}: Przystosowano rysowanie wykresów w bibliotece \verb|Matplotlib| by obsługiwały zarówno tryb prezentacji (Część Rzeczywista / Część Urojona), jak i format trygonometryczny (Moduł / Faza). Aby zwiększyć czytelność analizy widmowej wyeliminowano redundantne drugie połówki widma, co uwidoczniło zjawiska tuż za zerem (DC).
    \item \textbf{Skrypt testowy}: Przygotowano dedykowany skrypt porównujący zaimplementowane, autorskie metody na tle profesjonalnej biblioteki \verb|scipy.fft|. Skrypt automatycznie weryfikuje Błąd Średniokwadratowy (MSE) oraz zlicza stempel czasowy, weryfikując złożoność.
\end{itemize}

\section{Eksperymenty i wyniki}
\subsection{Eksperyment 1: Transformacje i ich transformacje odwrotne}
\subsubsection{Założenia i przebieg}
W ramach eksperymentu dla wszystkich trzech funkcji złożonych funkcji S1, S2 i S3 wykonano transformaty: odpowiednio Fouriera, Kosinusowa i Walsha-Hadamarda, a następnie wykonano transformaty odwrotne (w przypadku Kosinusowej i Walsha-Hadamarda transformata odwrotna jest taka sama jak zwykła transformata) w celu zbadania poprawności działania programu.

\subsubsection{Rezultaty}
\begin{figure}[H]
    \centering
    \begin{subfigure}[b]{0.32\textwidth}
        \centering
        \includegraphics[width=\textwidth]{ex1/S1.png}
        \caption{Sygnał S1}
    \end{subfigure}
    \hfill
    \begin{subfigure}[b]{0.32\textwidth}
        \centering
        \includegraphics[width=\textwidth]{ex1/S2.png}
        \caption{Sygnał S2}
    \end{subfigure}
    \hfill
    \begin{subfigure}[b]{0.32\textwidth}
        \centering
        \includegraphics[width=\textwidth]{ex1/S3.png}
        \caption{Sygnał S3}
    \end{subfigure}
    
    \caption{Oryginalne sygnały S1, S2 i S3 (przed transformacjami)}
\end{figure}

\begin{figure}[H]
    \centering
    \begin{subfigure}[b]{0.32\textwidth}
        \centering
        \includegraphics[width=\textwidth]{ex1/s1_DFT_1.png}
        \caption{Transformata: Rzeczywista/Urojona}
    \end{subfigure}
    \hfill
    \begin{subfigure}[b]{0.32\textwidth}
        \centering
        \includegraphics[width=\textwidth]{ex1/s1_DFT_2.png}
        \caption{Transformata: Moduł/Faza}
    \end{subfigure}
    \hfill
    \begin{subfigure}[b]{0.32\textwidth}
        \centering
        \includegraphics[width=\textwidth]{ex1/s1_DFT_odw.png}
        \caption{Odwrotna Transformata}
    \end{subfigure}
    
    \caption{Transformata Fouriera i odwrotna transformata Fouriera sygnału S1}
\end{figure}

\begin{figure}[H]
    \centering
    \begin{subfigure}[b]{0.48\textwidth}
        \centering
        \includegraphics[width=\textwidth]{ex1/s2_DCT.png}
        \caption{Transformata: Rzeczywista}
    \end{subfigure}
    \hfill
    \begin{subfigure}[b]{0.48\textwidth}
        \centering
        \includegraphics[width=\textwidth]{ex1/s2_DCT_odw.png}
        \caption{Odwrotna Transformata}
    \end{subfigure}
    
    \caption{Transformata Kosinusowa i jej odwrotność dla sygnału S2}
\end{figure}

\begin{figure}[H]
    \centering
    \begin{subfigure}[b]{0.48\textwidth}
        \centering
        \includegraphics[width=\textwidth]{ex1/s3_WHT.png}
        \caption{Transformata: Rzeczywista}
    \end{subfigure}
    \hfill
    \begin{subfigure}[b]{0.48\textwidth}
        \centering
        \includegraphics[width=\textwidth]{ex1/s3_WHT_odw.png}
        \caption{Odwrotna Transformata}
    \end{subfigure}
    
    \caption{Transformata Walsha-Hadamarda i jej odwrotność dla sygnału S3}
\end{figure}

 Wyniki odwrotnych transformacji pokrywają się z oryginalnymi wykresami. Oznacza to, że transformaty oraz ich transformaty odwrotne zostały prawidłowo zaprogramowane i działają zgodnie z założeniami.

\newpage
\subsection{Eksperyment 2: Porównanie złożoności czasowej algorytmów}
\subsubsection{Założenia i przebieg}
W ramach eksperymentu dla wszystkich trzech funkcji złożonych funkcji S1, S2 i S3 przeprowadzono porównanie transformaty Fouriera, Kosinusowej i Walsha-Hadamarda oraz ich wersjy szybkich.

\subsubsection{Rezultaty}
\begin{figure}[H]
    \centering
    \includegraphics[width=0.7\linewidth]{ex2/plot_S1.png}
    \caption{Wykres porównujący czas wykonywania poszczególnych transformat dla liczby próbek 512-8192 dla sygnału S1}
\end{figure}
\begin{figure}[H]
    \centering
    \includegraphics[width=0.7\linewidth]{ex2/plot_S2.png}
    \caption{Wykres porównujący czas wykonywania poszczególnych transformat dla liczby próbek 512-8192 dla sygnału S2}
\end{figure}
\begin{figure}[H]
    \centering
    \includegraphics[width=0.7\linewidth]{ex2/plot_S3.png}
    \caption{Wykres porównujący czas wykonywania poszczególnych transformat dla liczby próbek 512-8192 dla sygnału S3}
\end{figure}

Zauważamy, że zgodnie z przewidywaniami szybkie warianty transformat miały znacznie niższą złożoność czasową. Co ciekawe najwolniejszą transformatą okazała się transformata kosinusowa \verb|DCT| a najszybszą zazwyczaj był jej szybki odpowiednik \verb|FCT|, czasem ustępując jedynie szybkiej transformacie Walsha-Hadamarda \verb|FWHT|. Zauważamy marginalne różnice szybkiej transformaty Fouriera \verb|FFT| w dziedzinie częstotliwości względem dziedziny czasu.

\subsection{Eksperyment 3: Porównanie wyników transformat i szybkich transformat}
\subsubsection{Założenia i przebieg}
W ramach eksperymentu porównano wyniki transformat Fouriera, Kosinusowej i Walsha-Hadamarda z ich szybkimi odpowiednikami dla sygnału S3.

\subsubsection{Rezultaty}


\begin{figure}[H]
    \centering
    \begin{subfigure}[b]{0.48\textwidth}
        \centering
        \includegraphics[width=\textwidth]{ex3/tr_DFT_1.png}
        \caption{DFT - Rzeczywista/Urojona}
    \end{subfigure}
    \hfill
    \begin{subfigure}[b]{0.48\textwidth}
        \centering
        \includegraphics[width=\textwidth]{ex3/tr_DFT_2.png}
        \caption{DFT - Moduł/Faza}
    \end{subfigure}
    \vspace{0.5cm}
    \begin{subfigure}[b]{0.48\textwidth}
        \centering
        \includegraphics[width=\textwidth]{ex3/tr_FFT_1.png}
        \caption{FFT DIT - Rzeczywista/Urojona}
    \end{subfigure}
    \hfill
    \begin{subfigure}[b]{0.48\textwidth}
        \centering
        \includegraphics[width=\textwidth]{ex3/tr_FFT_2.png}
        \caption{FFT DIT - Moduł/Faza}
    \end{subfigure}
    \vspace{0.5cm}
    \begin{subfigure}[b]{0.48\textwidth}
        \centering
        \includegraphics[width=\textwidth]{ex3/tr_FFT_1_DIF.png}
        \caption{FFT DIF - Rzeczywista/Urojona}
    \end{subfigure}
    \hfill
    \begin{subfigure}[b]{0.48\textwidth}
        \centering
        \includegraphics[width=\textwidth]{ex3/tr_FFT_2_DIF.png}
        \caption{FFT DIF - Moduł/Faza}
    \end{subfigure}
    
    \caption{Porównanie Transformaty Fouriera DFT i Szybkiej Transformaty Kosinusowej FFT (w dziedzinie czasu DIT oraz dziedzinie częstotliwości DIF) dla sygnału S3}
\end{figure}
\begin{figure}[H]
    \centering
    \begin{subfigure}[b]{0.48\textwidth}
        \centering
        \includegraphics[width=\textwidth]{ex3/tr_DCT.png}
        \caption{DCT}
    \end{subfigure}
    \hfill
    \begin{subfigure}[b]{0.48\textwidth}
        \centering
        \includegraphics[width=\textwidth]{ex3/tr_FCT.png}
        \caption{FCT}
    \end{subfigure}
    
    \caption{Porównanie Transformaty Kosinusowej DCT i Szybkiej Transformaty Kosinusowej FCT dla sygnału S3}
\end{figure}
\begin{figure}[H]
    \centering
    \begin{subfigure}[b]{0.48\textwidth}
        \centering
        \includegraphics[width=\textwidth]{ex3/tr_WHT.png}
        \caption{WHT}
    \end{subfigure}
    \hfill
    \begin{subfigure}[b]{0.48\textwidth}
        \centering
        \includegraphics[width=\textwidth]{ex3/tr_FWHT.png}
        \caption{FWHT}
    \end{subfigure}
    
    \caption{Porównanie Transformaty Walsha-Hadamarda WHT i Szybkiej Transformaty Walsha-Hadamarda FWHT dla sygnału S3}
\end{figure}

Na podstawie powyższych wyników zauważamy, że transformaty dają takie same wyniki co ich szybkie odpowiedniki, jak również takie same wyniki dla Szybkich Transformat Fouriera w różnych dziedzinach. Jest to zgodne z założeniami i potwierdza poprawność zaimplementowanych transformat.

\newpage
\section{Wnioski}
Zaimplementowane transformacje Fouriera, Kosinusowe oraz Walsha-Hadamarda pozwalają na sprawną dekompozycję sygnałów z dziedziny czasu na wybrane przez nas częstotliwości i ortogonalne składowe bazowe.\\
Możliwość wizualizacji za pomocą trybu Moduł/Faza w doskonały sposób ukazuje rzeczywisty rozkład amplitudowy - pomijając zjawiska przesunięć oraz ułatwia odczyt składowych sygnałów złożonych.\\
Implementacja klasycznych form transformat udowadnia, dlaczego tak rzadko stosuje się je we współczesnych narzędziach - ich czas rośnie w czasie wielomianowym ($N^2$), co kompletnie dyskwalifikuje je w ujęciach Real-Time (na bieżąco). Szybkie transformaty jednak wprowadzają ograniczenie - liczba próbek poddawana przekształceniu w danym cyklu musi być potęgą dwójki.\\
Użycie bibliotek takich jak wbudowane do \verb|scipy| i \verb|numpy| weryfikuje naszą pewność co do działania aplikacji.


\end{document}